% 03-decomposition.tex — OpenBench Software Design Document. Written from .scratch/design-doc/brief.md on 2026-10-05.
\section{Problem Decomposition}\label{sec:decomp}

The problem was broken down in two passes. The first pass turned the twelve features of the SRS, plus the concerns that cut across them, into responsibilities small enough for one module to own. The second pass took the four parts that resist a clean split (three allocation policies, a Request whose parts resolve at different speeds, credits, and timed work) and decided how each is divided. Section~\ref{sec:decomp-resp} gives the result of the first pass, Section~\ref{sec:decomp-sub} the reasoning of the second.

\subsection{Responsibilities}\label{sec:decomp-resp}

Every responsibility has exactly one owning module, the one that holds its data and enforces its rules; collaborators are the modules it calls or is called by. Where two modules share a responsibility, the first named owns the decision and the second carries it out. Every one of the 41 Release~1 functional requirements appears in at least one row, and Appendix~\ref{app:trace} gives the reverse mapping. The module specifications of Section~\ref{sec:modules} cite these identifiers.

{\small\renewcommand{\arraystretch}{1.08}\setlength{\tabcolsep}{4pt}
\begin{longtable}{@{}L{0.75cm}L{0.31\textwidth}L{0.29\textwidth}L{0.13\textwidth}L{\dimexpr0.27\textwidth-0.75cm-8\tabcolsep\relax}@{}}
\toprule
\textbf{Id} & \textbf{Responsibility} & \textbf{SRS identifiers} & \textbf{Owner} & \textbf{Collaborators} \\
\midrule
\endfirsthead
\multicolumn{5}{@{}l}{\small\itshape Responsibilities (continued)}\\ \toprule
\textbf{Id} & \textbf{Responsibility} & \textbf{SRS identifiers} & \textbf{Owner} & \textbf{Collaborators} \\
\midrule
\endhead
\midrule \multicolumn{5}{@{}r}{\small\itshape continues on next page}\\
\endfoot
\bottomrule
\endlastfoot
\multicolumn{5}{@{}l}{\itshape Inventory and profiles}\\
\rid{1} & Define, change, and retire Resource Types as data, with one Allocation Class and typed attributes; protect the built-in type Person & \req{Inventory.Type.Define}, \req{QA-extensible} & \module{Inventory} & \module{Members} \\
\rid{2} & Keep Resources, their units, and Pool Capacity; refuse a retirement or reduction below future usage, naming the earliest hour & \req{Inventory.Resource.Add} & \module{Inventory} & \module{Allocation} (\code{ReductionGuard}) \\
\rid{3} & Show who holds a Resource in each part of a window; filter Resources by type, attribute, and free for a window & \req{Inventory.View.Holder}, \req{UI-inventory} & \module{Allocation} & \module{Inventory}, Web client \\
\rid{4} & Keep Members, Teams with exactly one Manager, roles, the skill list, each Member's skills and availability; answer who is available with which skills & \req{Admin.Member.Create}, \req{Profile.Skills}, \req{Profile.Availability}, \req{SE-data} & \module{Members} & \code{security} \\
\addlinespace
\multicolumn{5}{@{}l}{\itshape Suggestions, drafting, and submission}\\
\rid{5} & Create Suggestions; accept one into a new or existing line, or decline it with a reason & \req{Suggest.Create}, \req{Suggest.Accept.Existing} & \module{Requests} & \module{Members}, \module{Notifications} \\
\rid{6} & Edit Drafts: add, edit, split, and remove lines; mark Required or Optional and Priority; refuse past windows and retired types & \req{Draft.Create}, \req{Draft.Line.Add}, \req{Draft.Line.Required}, \req{UI-draft} & \module{Requests} & \module{Inventory}, \module{Settings} \\
\rid{7} & Price a line (Quote) and a Release (Refund fraction) with the pure class \code{Pricing}, from the line, the settings, and the lead time & \req{Ledger.Price.Formula}, \req{Draft.Quote.Line}, \req{Metered.Line.Form}, \req{Release.Refund} & \module{Ledger} & \module{Settings}, \module{Requests} \\
\rid{8} & Check affordability and submit: refuse a shortfall in Required lines, or any priced Required line while the Balance is negative; lock every Quote; record the submission time & \req{Draft.Quote.Line}, \req{Draft.Submit.Lock}, \req{Release.Loan.Prorate} & \module{Requests} & \module{Ledger} \\
\addlinespace
\multicolumn{5}{@{}l}{\itshape Allocation per class}\\
\rid{9} & Decide Matched eligibility (skills, availability, other Team, active, no overlapping Loan) and build the Shortlist least-lent first; form Loans, one per Person at a time & \req{Matched.Shortlist.Build}, \req{Profile.Availability}, \req{Matched.Loan.Form}, \req{PE-shortlist}, \req{UI-shortlist} & \module{Allocation} & \module{Members}, \module{Requests} \\
\rid{10} & Route one approval per selected Candidate to the lending Manager; record each decision; on refusal release Holds and ask for a replacement, or refuse or drop the line & \req{Matched.Approve.Route}, \req{Matched.Refuse.Replace}, \req{UI-shortlist} & \module{Allocation}, \module{Requests} & \module{Members}, \module{Ledger}, \module{Notifications} \\
\rid{11} & Detect Exclusive conflicts against Held and Allocated occurrences and pick a free instance; never double-allocate & \req{Exclusive.Conflict}, \req{QA-correct} & \module{Allocation} & \module{Inventory} \\
\rid{12} & Check a Metered line hour by hour against Pool Capacity and the Team's Quota; never oversell & \req{Metered.Line.Form}, \req{Metered.Pool.Check}, \req{QA-correct} & \module{Allocation} & \module{Inventory}, \module{Settings} \\
\rid{13} & Keep the Waitlist: join, order (Priority, then submission time, Bumped at the head), positions, and expiry at window start & \req{Exclusive.Waitlist.Join}, \req{Metered.Pool.Check} & \module{Allocation} & \module{Requests} \\
\rid{14} & Plan and apply Bumps for Priority lines outside the bump cutoff; refund the Bumped line in full and put it at the head of the Waitlist & \req{Exclusive.Priority.Bump}, \req{Exclusive.Bump.Effect} & \module{Allocation}, \module{Requests} & \module{Ledger}, \module{Settings}, \module{Notifications} \\
\addlinespace
\multicolumn{5}{@{}l}{\itshape Request lifecycle}\\
\rid{15} & Run the Request state machine; evaluate every line on submission and choose Allocated, Pending, or Waiting & \req{Lifecycle.States}, \req{Lifecycle.Evaluate} & \module{Requests} & \module{Allocation} \\
\rid{16} & Place Holds with held credits when a Request enters Pending; release them when it leaves Pending; place none after the window lapses & \req{Lifecycle.Hold.Place} & \module{Requests} & \module{Allocation}, \module{Ledger} \\
\rid{17} & Allocate a Request in one transaction: commit Holds, post Charges and Earnings, form Loans, drop Optional lines with a reason & \req{Lifecycle.Allocate}, \req{Matched.Loan.Form}, \req{QA-atomic}, \req{CO-tx} & \module{Requests} & \module{Allocation}, \module{Ledger}, \module{Notifications} \\
\rid{18} & Start the Objection Window on first entry to Pending; on lapse release Holds, move to Waiting, and escalate to the Capacity Planner & \req{Lifecycle.Window.Start}, \req{Matched.Escalate}, \req{Planner.Settings} & \module{Requests} & \module{Allocation}, \module{Ledger}, \module{Jobs} \\
\rid{19} & Re-evaluate Waiting Requests when capacity frees or an approval arrives; allocate one only when every Required line fits at once & \req{Lifecycle.Waiting.Reevaluate}, \req{Exclusive.Waitlist.Join} & \module{Requests} & \module{Allocation}, \module{Ledger} \\
\addlinespace
\multicolumn{5}{@{}l}{\itshape Release and Ledger}\\
\rid{20} & Release a line in whole or in part and end a Loan early; post the Refund or the Earning correction; offer the freed capacity & \req{Release.Line}, \req{Release.Refund}, \req{Release.Loan.Prorate} & \module{Requests} & \module{Allocation}, \module{Ledger} \\
\rid{21} & Keep Accounts (granted, earned, held) and append-only Transactions with the spend-order split; post Earnings; produce Statements & \req{Ledger.Account}, \req{Ledger.Transaction.Kinds}, \req{Ledger.Earning.Post}, \req{QA-audit} & \module{Ledger} & \module{Settings}, \module{Members} \\
\rid{22} & Post planned and mid-Period Grants; roll the Period over, expiring granted credits and keeping earned ones & \req{Ledger.Grant.Period}, \req{PE-jobs} & \module{Ledger} & \module{Jobs}, \module{Settings} \\
\addlinespace
\multicolumn{5}{@{}l}{\itshape Reporting and settings}\\
\rid{23} & Show utilization, idle capacity, waiting lines, credits, and person-days, scoped to the viewer's roles & \req{Dashboard.Utilization}, \req{Dashboard.Credits}, \req{UI-inventory} & \module{Dashboard} & \module{Members} \\
\rid{24} & Keep every Capacity Planner setting as data, applied to future Quotes and events, with a change log & \req{Planner.Settings}, \req{Admin.Member.Create} & \module{Settings} & \module{Inventory}, \module{Members} \\
\addlinespace
\multicolumn{5}{@{}l}{\itshape Crosscutting}\\
\rid{25} & Authenticate Members by username and password; enforce roles on every action, from the client or the API alike & \req{Admin.Member.Create}, \req{SE-roles}, \req{SE-password}, \req{CO-auth}, \req{CO-api} & \code{security}, \module{Members} & every module \\
\rid{26} & Scope every read to the caller's own record, Team, owned Resources, or role & \req{SE-visibility}, \req{Profile.Availability}, \req{Ledger.Account}, \req{Dashboard.Credits} & every module & \code{security} \\
\rid{27} & Record actor, change, and time for every Request and line state change and every setting change & \req{Lifecycle.States}, \req{Exclusive.Bump.Effect}, \req{QA-audit}, \req{QA-retention} & \module{Requests}, \module{Settings} & \module{Ledger} \\
\rid{28} & Run timed work: lapse Objection Windows, expire waiting lines, roll Periods over, from one injected \code{Clock} & \req{Lifecycle.Window.Start}, \req{Ledger.Grant.Period}, \req{PE-jobs}, \req{CO-jobs} & \module{Jobs} & \module{Requests}, \module{Ledger}, \module{Settings} \\
\rid{29} & Notify Members in-app of approval requests, outcomes, Bumps, Loans, escalations, and declined Suggestions & \req{CI-inapp}, \req{Lifecycle.Allocate}, \req{Matched.Loan.Form}, \req{Suggest.Accept.Existing} & \module{Notifications} & \module{Requests}, \module{Allocation} \\
\rid{30} & Present the role-based screens over the REST API, with Quotes fetched from the server & \req{UI-draft}, \req{UI-shortlist}, \req{UI-inventory}, \req{QA-usable}, \req{QA-available} & Web client & every controller \\
\addlinespace
\multicolumn{5}{@{}l}{\itshape Stretch, designed as extension points (Section~\ref{sec:evo})}\\
\rid{31} & Service tiers with a premium on Metered lines, and Bumps of lower-tier units by \code{Bumpable} on \code{MeteredAllocator} & \req{Metered.Tier.Define}, \req{Metered.Tier.Bump} & \module{Allocation} & \module{Settings}, \module{Ledger} \\
\rid{32} & Record units actually used per Metered allocation and compare requested with used & \req{Metered.Usage.Record}, \req{Metered.Usage.Compare}, \req{Dashboard.Usage} & \module{Allocation}, \module{Dashboard} & \module{Inventory} \\
\rid{33} & Add a demand surcharge multiplier inside \code{Pricing.quote} that rises as the requested window fills & \req{Ledger.Surcharge} & \module{Ledger} & \module{Settings}; \module{Requests} passes the fill from \module{Allocation} \\
\rid{34} & Deliver notifications by email as well as in-app, through a second \code{Delivery}, \code{EmailDelivery} & \req{CI-inapp} (Stretch part) & \module{Notifications} & \module{Members} \\
\end{longtable}
}

\subsection{Subproblems and How Each Was Broken Down}\label{sec:decomp-sub}

\subsubsection{Three allocation policies behind one interface}

The three Allocation Classes answer the same question, ``can this line have its capacity, and on what terms?'', in different ways. An Exclusive line checks occurrences for overlap and may Bump; a Metered line sums units hour by hour; a Matched line waits for a lending Manager's approval. The lifecycle should not care which. The breakdown separated what the lifecycle needs from every class from what only some classes do. Every class must evaluate a line, hold it, commit it, release it, and cancel it, so those five operations form \code{Allocator}. The rest became three optional capabilities that \module{Requests} asks the \code{AllocatorRegistry} for and skips when a class lacks them (\adr{0007}):

{\small\renewcommand{\arraystretch}{1.05}
\begin{tabularx}{\textwidth}{@{}L{3.1cm}L{4.4cm}YYY@{}}
\toprule
\textbf{Allocator} & \textbf{\code{Allocator} (five operations)} & \textbf{\code{Waitlistable}} & \textbf{\code{Bumpable}} & \textbf{\code{ApprovalDriven}} \\
\midrule
\code{ExclusiveAllocator} & all; \code{hold} writes Held occurrences & yes & yes & no \\
\code{MeteredAllocator} & all; \code{hold} writes Held units & yes & no (Stretch: yes) & no \\
\code{MatchedAllocator} & all; \code{hold} is a no-op, \code{commit} forms Loans & no & no & yes \\
\bottomrule
\end{tabularx}
}

The parts most likely to be wrong are pure functions inside the allocators, testable without a database: conflict detection over windows and weekly recurrences, the hourly sum against Pool Capacity and Quota, and the Shortlist filter and order. \emph{Rejected:} one wide interface with every operation, which forces Matched to fake a Waitlist and every new class to implement methods that mean nothing for it; and three unrelated services that \module{Requests} picks by class, which would make every new class edit the orchestrator, against \req{CO-plugin} and \req{QA-extensible}.

\subsubsection{A Request whose parts resolve at different speeds}

Exclusive and Metered lines resolve at once against the database; a Matched line waits days for a person. Allocating the room before the approvals arrive charges for a room that may go unused; waiting for the approvals loses the room. The breakdown has three parts. First, one orchestrator: \code{RequestLifecycleService} in \module{Requests} owns the state machine and is the only code that moves a Request between states. Second, Holds bridge the gap (\adr{0001}): when a Request enters Pending, every satisfiable Exclusive and Metered line is Held, blocking capacity without a Charge, and its Quote counts as held credits, all bounded by the Objection Window. Third, one transaction per transition: submission, an approval decision, a lapse, and a Release (with the waitlist fill it triggers) are each one \code{@Transactional} method that calls the allocators and \module{Ledger} directly, so a Request is never seen half allocated. Waiting is the single recovery state for refusals, lapses, and Bumps, which keeps the state machine to seven states. \emph{Rejected:} a saga of modules reacting to each other's events, with a compensating step for each failure, which exposes half-allocated Requests and needs compensation logic for every transition; and best-effort allocation that charges whatever fits now, which charges for GPU hours nobody can use when the people are refused.

\subsubsection{Credits as pure pricing plus append-only posting}

Credits raise two different questions: what a line costs, and what happened to an Account. They were split. \emph{Pricing} is a pure class, \code{Pricing}, in \module{Ledger}: \code{quote} sums quantity times Unit Price times Price Band multiplier over every hour, times the priority premium for a Priority line, and \code{refund} applies the Refund Schedule to a lead time. It reads no database and is used both while drafting and at submission, and the Quote it returns is locked on the Request Line, so the Charge always equals the Quote shown. \emph{Posting} is \code{LedgerService}: each operation updates the Team's \code{account} row (granted, earned, held) under that row's lock, in the caller's transaction, and every Grant, Charge, Earning, or Refund also appends a Transaction row. Affordability is a check of the Available Balance on the locked row; the spend order (granted first, then earned) is recorded on each Charge so a Refund can return credits in the same proportions (\adr{0002}). Matched lines never reach pricing at all (\adr{0004}). \emph{Rejected:} computing the Charge at commit time from the settings then in force, which lets a Charge differ from the Quote the Manager accepted; and deriving every balance by summing Transactions, which still needs a lock to be safe and grows slower with history.

\subsubsection{Timed work as polling in one process}

Four things happen because time passes: an Objection Window lapses, a Hold expires with it, a waiting line's window start passes, and a Period ends. Each was reduced to a row whose time has come: a Pending or Waiting Request whose \code{window\_ends\_at} is past with an approval not yet escalated, a waiting line whose start is past, and a Period whose end is past. \code{DueWorkJob} polls for the first two every 60~seconds and \code{PeriodRolloverJob} handles the third at 00:05 on the boundary day and again at start-up. Each due item is handled in its own transaction by an application-service method of the owning module (\code{onWindowLapse}, \code{expireWaitingLines}, \code{rollover}) that reuses the steps user actions use, such as releasing Holds and moving a Request to Waiting, with the system recorded as the actor. Because each query selects only items still in the due state, a second run finds nothing to do, and a restart loses nothing. A \code{java.time.Clock} bean is injected wherever time is read, so tests move time instead of waiting. \emph{Rejected:} scheduling one timer per due item, which loses timers on restart unless a job store such as Quartz is added; operating-system cron, which \req{CO-jobs} rules out; and scheduling inside the database, which would move lifecycle logic out of the code that owns it.
